% ─────────────────────────────────────────────────────────────────────────────
%  Sección 3.9  Módulo de Ejercicios y Retroalimentación
%  Insertar al final del Capítulo 3: Análisis, después de la sección 3.8
%  (Criterios de selección de dataset).
%  CITES: Reemplazar sep_tercer_grado, sep_cuarto_grado y sep_planes_primaria
%  con las claves exactas de su archivo .bib cuando las tengan.
% ─────────────────────────────────────────────────────────────────────────────

\subsection{Módulo de Ejercicios y Retroalimentación}

El módulo de ejercicios y retroalimentación es el componente pedagógico
central del sistema, encargado de transformar los resultados de los
análisis ortográfico y caligráfico en acciones correctivas concretas.
Su propósito es proporcionar al alumno material de práctica focalizado
en sus áreas de oportunidad específicas y al tutor las herramientas para
supervisar y complementar dicho proceso. A diferencia de los módulos de
captura y análisis, cuya naturaleza es predominantemente técnica, este
módulo establece el puente entre el procesamiento computacional y el
contexto pedagógico en el que opera la aplicación.

% ─────────────────────────────────────────────────────────
\subsubsection{Justificación pedagógica de los ejercicios}

La selección y diseño de los ejercicios dentro de la plataforma se
fundamentan en los requerimientos del plan de estudios de la Secretaría
de Educación Pública (SEP) para tercer y cuarto grado de primaria
\cite{sep2011guia, sep2024saberes4, mejoredu2022indicadores}, etapa
en la cual el alumno transita de la decodificación básica hacia la
consolidación de la ortografía convencional y la fluidez del trazo
manuscrito. En este periodo, los errores no corregidos tienden a fijarse
como hábitos, lo que justifica la intervención temprana y personalizada
que el sistema propone.

Los ejercicios se categorizan en dos grandes rubros, en correspondencia
directa con las métricas evaluadas por el sistema:

\paragraph{Ejercicios de refuerzo caligráfico}

Orientados a mejorar las métricas de trazo extraídas por el módulo OCR:
alineación, tamaño de letra, espaciado e inclinación. Se contemplan dos
tipos:

\begin{enumerate}
    \item \textbf{Copia de modelos (pangramas y frases cortas):} Se
    utilizan frases que contienen todas las letras del alfabeto
    —pangramas— y oraciones adaptadas al vocabulario infantil. La
    práctica de este tipo de textos obliga al alumno a trabajar la
    proporción de tamaño y el espaciado correcto entre caracteres de
    distinta frecuencia de uso, dado que los pangramas incluyen letras
    poco comunes que requieren mayor atención en el trazo.

    \item \textbf{Trazos de direccionalidad:} Ejercicios enfocados en
    grupos de letras que comparten patrones motores similares, como
    \textit{a, c, d, g, q} o \textit{b, p, d, q}. Su objetivo es
    corregir la inclinación del trazo y reducir la confusión espacial
    derivada de la similitud morfológica entre estos caracteres, un
    fenómeno frecuente en alumnos de este nivel educativo.
\end{enumerate}

\paragraph{Ejercicios de corrección ortográfica}

Generados dinámicamente a partir de los errores detectados por el módulo
NLP y asignados según el tipo de falla identificada. Se contemplan tres
modalidades:

\begin{enumerate}
    \item \textbf{Uso de grafías conflictivas:} Ejercicios de completar
    palabras enfocados en las reglas ortográficas de mayor dificultad en
    este nivel educativo: uso de \textit{b/v}, \textit{c/s/z},
    \textit{g/j} y el uso de la \textit{h}. Cada ejercicio presenta al
    alumno el contexto de la palabra con la grafía omitida, de modo que
    la decisión ortográfica requiera razonamiento y no memorización
    mecánica.

    \item \textbf{Acentuación:} Ejercicios de identificación y
    colocación de tildes, con énfasis en la clasificación de palabras en
    agudas, graves y esdrújulas. El sistema selecciona palabras del mismo
    campo semántico que el texto original del alumno para mantener la
    coherencia con el ejercicio que motivó el análisis.

    \item \textbf{Reescritura focalizada:} Cuando el módulo NLP detecta
    una palabra mal escrita en el texto libre del alumno, el sistema
    genera automáticamente un ejercicio de repetición espaciada en el
    que el alumno debe reescribir correctamente esa palabra específica
    dentro de un nuevo contexto oracional. Esta modalidad es la de mayor
    personalización, pues el contenido del ejercicio se construye
    directamente a partir del error individual del alumno.
\end{enumerate}

% ─────────────────────────────────────────────────────────
\subsubsection{Estructura del catálogo de ejercicios}

El catálogo de ejercicios disponibles en la plataforma se almacena en
la tabla \texttt{Ejercicios} de la base de datos. Cada registro del
catálogo contiene los atributos necesarios para que el sistema pueda
tanto presentarlo al alumno como procesarlo automáticamente para la
generación de variantes. La Tabla~\ref{tab:catalogo-ejercicios} describe
los atributos del catálogo y su propósito dentro del módulo.

\begin{table}[H]
\centering
\caption{Atributos del catálogo de ejercicios del sistema}
\label{tab:catalogo-ejercicios}
\begin{tabular}{|l|l|p{7cm}|}
\hline
\textbf{Atributo} & \textbf{Tipo} & \textbf{Descripción} \\
\hline
\texttt{id\_ejercicio}  & Entero        & Identificador único del ejercicio en el catálogo. \\
\texttt{titulo}         & VARCHAR(150)  & Nombre descriptivo del ejercicio visible para el tutor. \\
\texttt{descripcion}    & VARCHAR(MAX)  & Instrucciones del ejercicio para el alumno. \\
\texttt{tipo}           & VARCHAR(30)   & Categoría del ejercicio: \texttt{caligrafico} u \texttt{ortografico}. \\
\texttt{contenido\_base}& VARCHAR(MAX)  & Texto base, pangrama, frase o patrón de letra sobre el que se construye el ejercicio. \\
\texttt{id\_estatus}    & Entero        & Estado del ejercicio en el catálogo (activo/inactivo). \\
\hline
\end{tabular}
\end{table}

El campo \texttt{tipo} actúa como discriminador para que el sistema
pueda filtrar y asignar ejercicios de forma coherente con el tipo de
error detectado: errores caligráficos activan ejercicios de tipo
\texttt{caligrafico}, mientras que errores ortográficos activan
ejercicios de tipo \texttt{ortografico}. Esta separación garantiza que
la retroalimentación entregada al alumno sea siempre pertinente al área
de oportunidad identificada en su intento.

La relación entre el catálogo y el tutor se gestiona a través de la
tabla \texttt{Ejercicios\_Tutor}, que representa la activación de un
ejercicio del catálogo por parte de un tutor específico y contiene su
propio ciclo de vida independiente, conforme al diseño descrito en la Sección~4.5.
